\chapter{Diskussion}
\label{chap:diskussion}

Die Ergebnisse zeigen, dass ungelabelte Bilder das Retrieval auf ungesehenen Klassen auch gegenüber einer supervised Referenz (\UEBL) verbessern können, allerdings nur unter bestimmten Bedingungen und breit nur mit Label Propagation. Dieses Kapitel fragt deshalb weniger, welches Verfahren auf welchem Datensatz gewinnt, als unter welchen Bedingungen sich ungelabelte Bilder lohnen. Die Forschungsfragen und die Hypothesen werden dort beantwortet, wo die jeweilige Evidenz besprochen wird.



\section{Wann sich ungelabelte Bilder lohnen}
\label{sec:d:nutzen}

\subsection{Nutzen gegenüber einer abgestimmten Referenz}

Die erste Forschungsfrage (F1) lautete, ob ungelabelte Bilder die Retrieval-Leistung auf ungesehenen Testklassen gegenüber einer supervised Referenz verbessern. Die Antwort ist ein bedingtes Ja. Belastbar ist der Gewinn vor allem bei den Budgets, bei denen die Referenz und die Mechanismen abgestimmt wurden, auf Cars196 bei $k=16$ für \ac{LP} und auf In-Shop bei $k=2$ für \ac{LP} und Hoffer, die dort in keinem gepaarten Vergleich unter der Referenz liegen (Abschnitt~\ref{sec:e:phase1}). Den größten Gewinn erreicht \ac{LP} auf Cars196 bei vier gelabelten Bildern je Klasse. Im Mittel über die Verlustfunktionen ist die Richtung dieses Gewinns plausibel, denn um ihn aufzuheben, müsste eine für vier Labels je Klasse neu abgestimmte Referenz die bei $k=16$ abgestimmte Referenz mit viermal so vielen Labels übertreffen. Die Größe des Gewinns ist dagegen weder eine obere noch eine untere Schranke. Eine besser regularisierte Referenz könnte ihn verkleinern, ein für $k=4$ abgestimmtes \ac{LP} ihn vergrößern.%


In Trainingslabels ausgedrückt, bei derselben vollständig gelabelten Validierung in allen Budgets (Abschnitt~\ref{sec:budgets}), erreicht \ac{LP} auf Cars196 mit vier und mit acht gelabelten Bildern je Klasse im Mittel etwa die Leistung der Referenz mit der vierfachen Zahl, allerdings nur in ungepaarten Vergleichen innerhalb der gemessenen Budgets. Hypothese~H1 ist damit teilweise gestützt. Für \ac{LP} gilt sie auf Cars196 und Semi-iNat-2021 bis $k=16$ und auf In-Shop bei $k=2$ ohne einen negativen gepaarten Vergleich, für Hoffer, \ac{SLADE} und \ac{LRML} nur auf einzelnen Datensätzen (Abschnitt~\ref{sec:e:phase1}). Auf \ac{CUB} ist sie nur für das Ensemble mit wenigen Zehnteln eines Prozentpunkts gestützt.%


Mit früheren Arbeiten zum \ac{SSDML} lassen sich diese Gewinne nur bedingt vergleichen. Auf \ac{CUB} steigt in einer davon mit fünf Labels je Klasse die \ac{P@1}, dort als Recall@1 berichtet, von \qty{31.7}{\percent} vor dem Training auf \qty{44.8}{\percent}, aber nur um \num{2.2} bis \num{3.4} Prozentpunkte über supervised Baselines mit denselben Labels~\cite{dutta_semi-supervised_2021}. Label Propagation mit proxybasierten Ankern wird auf einer eigenen Aufteilung der Klassen gegen supervised Verfahren mit Standard-Hyperparametern verglichen und liegt in der \ac{NMI} auf Cars196 um 2,6 und auf \ac{CUB} um 4,4 Punkte über dem besten von ihnen~\cite{zhuang_deep_2021}. Die gemischte Label Propagation der Folgearbeit liegt dort um 2,6 und 4,5 Punkte darüber~\cite{zhuang_deep_2023}. Alle drei Arbeiten trainieren ein vortrainiertes \ac{CNN} mit und berichten keine \ac{MAP@R}. Diese berichten \ac{SLADE} und \ac{STML}, beide aber mit allen Labels und zusätzlichen externen Bildern~\cite{duan_slade_2021, kim_self-taught_2022}. \ac{SLADE} übertrifft den eigenen Teacher mit derselben Verlustfunktion auf \ac{CUB} um \num{4.3} bis \num{4.7} und auf Cars196 um \num{3.7} bis \num{5.4} Prozentpunkte~\cite{duan_slade_2021}. \ac{STML} liegt mit einem auf ImageNet vortrainierten ResNet50 und 512 Dimensionen auf \ac{CUB} um \num{11.4} und auf Cars196 um \num{10.1} Prozentpunkte über dem supervised Modell mit \ac{MS}, von dem es ausgeht~\cite{kim_self-taught_2022}. In dieser Arbeit entfällt auf Cars196 und In-Shop ebenfalls der größere Teil des Sprungs von den unveränderten Merkmalen auf das supervised Training, auf Cars196 bei $k=4$ gut zwölf Prozentpunkte gegenüber weiteren \num{7.8} durch \ac{LP} (Tabelle~\ref{tab:e:referenz}, Abbildung~\ref{fig:e:gewinne}). Ein Gewinn gegenüber unveränderten Merkmalen oder nicht abgestimmten Baselines sagt damit wenig über den Nutzen ungelabelter Bilder aus. Auch in der Klassifikation schrumpften die berichteten Gewinne, sobald die supervised Baseline gleich intensiv abgestimmt wurde~\cite{oliver_realistic_2018} (Abschnitt~\ref{sec:ssl:evaluation}).

\subsection{Messbarer und praktisch bedeutsamer Gewinn}

Ob ein Gewinn messbar ist und ob er lohnenswert ist, ist nicht dasselbe. Auf Semi-iNat-2021 ist der Gewinn von \ac{LP} bei zwei Labels je Art kein Zufall der Seeds, entspricht aber nur dem Schritt der Referenz von zwei auf vier Labels je Art. Allerdings brächten dort auch zusätzliche echte Labels vergleichsweise wenig (Abschnitt~\ref{sec:e:ausgangslage}). Auf \ac{CUB} übertreffen schon die unveränderten DINOv2-Merkmale die meisten trainierten Einzelmodelle (Abschnitt~\ref{sec:e:ausgangslage}) und kein Mechanismus hat diesen Rückstand in einen nennenswerten Gewinn verwandelt. Auch beim Fine-Tuning des ganzen DINOv2 war \ac{CUB} unter mehreren Retrieval-Datensätzen der einzige, auf dem das unveränderte Modell kaum zurücklag~\cite{berton_all_2025}. Ebenso war \ac{CUB} unter den drei Datensätzen von ooDML der einzige, auf dem die unveränderten Merkmale eines vortrainierten ResNet (BiT) in der Recall@1 vor den dort trainierten \ac{DML}-Modellen lagen~\cite{milbich_characterizing_2021}.

Wie zu erwarten, fallen die Gewinne dort groß aus, wo auch zusätzliche echte Labels die Referenz noch deutlich verbessern. Der stärkste Gewinn wird auf Cars196 gemessen und der kleinste Gewinn auf \ac{CUB}. Das ist plausibel für Mechanismen, die Pseudo-Labels wie echte Labels verwenden, denn im günstigsten Fall wirkt jedes ihrer Pseudo-Labels wie ein echtes Label für dasselbe Bild. Den verbleibenden Spielraum zeigt die Referenz mit allen Labels der Trainingsklassen, die auf \ac{CUB} und Semi-iNat-2021 kaum, auf Cars196 und In-Shop dagegen noch über dem größten Budget liegt (Abschnitt~\ref{sec:e:ausgangslage}). Für einen Term wie den von \ac{LRML} ist das keine feste Grenze, denn auf Cars196 liegt \ac{LRML} auch mit allen Labels und ohne ungelabelte Bilder über der \UEBL (Abschnitt~\ref{sec:e:extern}). Mit vier Datensätzen ist dieser Zusammenhang eine Beobachtung. Der verbleibende Spielraum ist keine hinreichende Bedingung für einen Gewinn. Auf Semi-iNat-2021 gewinnt die Referenz von $k=2$ bis zu allen Labels noch rund drei Prozentpunkte, \ac{LP} bei $k=2$ aber nur \num{0.6}. In der In-Shop-Serie mit vollständig gelabelten Klassen gewinnt die Referenz mit mehr gelabelten Klassen ebenfalls deutlich, während \ac{LP} trotzdem verliert (Abschnitt~\ref{sec:d:daten}).
 

Die Auswahl der Vortrainingsdaten erklärt dieses Muster nicht allein (Abschnitt~\ref{sec:vit}). In dieser Arbeit bleibt trotzdem auf Cars196 der größte und auf Semi-iNat-2021 ein kleiner Spielraum. Für \ac{CUB} und Cars196 kommt als weitere Erklärung der Abstand zwischen Trainings- und Testklassen in Frage (Abschnitt~\ref{sec:modellauswahl}). Gemessen am Spektrum der ooDML-Aufteilungen liegt er in der üblichen Aufteilung von Cars196 nahe am Wert einer zufälligen Aufteilung der Bilder derselben Klassen, bei \ac{CUB} deutlich darüber. Für diese Arbeit folgt daraus nur die Hypothese, dass sich das auf den Trainingsklassen Gelernte auf Cars196 leichter auf die Testklassen überträgt als auf \ac{CUB}, denn der Abstand wurde mit ResNet50- und nicht mit DINOv2-Merkmalen gemessen \cite{milbich_characterizing_2021}. Auf Cars196 finden die unveränderten Merkmale meist ein Bild desselben Modells als ersten Treffer, ordnen die übrigen Bilder der Klasse aber nicht zusammen (Abschnitt~\ref{sec:e:ausgangslage}). Nach dem vollständigen Fine-Tuning eines auf ImageNet-21k vortrainierten ViT-S/16 mit Proxy-Anchor und allen Labels erreicht ein Modell auf Cars196 eine \ac{P@1}, als Recall@1 berichtet, von \qty{86.7}{\percent}, aber nur eine \ac{MAP@R} von \qty{28.7}{\percent}, auf \ac{CUB} bei ähnlicher \ac{P@1} dagegen \qty{51.1}{\percent}~\cite{ren_learning_2024}. Eine feste Eigenschaft von Cars196 ist das aber nicht, denn nach dem Fine-Tuning des ganzen DINOv2-B erreichen dort zehn von zwölf Verlustfunktionen bei Batchgröße 256 eine \ac{MAP@R} von rund \qty{54}{\percent} bis \qty{57}{\percent}, ähnlich wie auf \ac{CUB}~\cite{berton_all_2025}. Auch der lineare Klassifikationsvergleich aus Abschnitt~\ref{sec:vit} spricht dafür, dass die Merkmale die Klassen unterscheiden können. Im Beispiel aus Abbildung~\ref{fig:e:retrieval} folgt die Ähnlichkeit auf Cars196 aber auch der Farbe und der Größenklasse, während eine Klasse dasselbe Modell in jeder Farbe umfasst. Gerade die Lücke zwischen dem nächsten Nachbarn und der ganzen Klasse ist der Raum, den ein Head und mit ihm die Pseudo-Labels füllen können. Für \ac{SSL} mit vortrainierten Modellen eignen sich in der Klassifikation deshalb vor allem Aufgaben, bei denen eingefrorene Merkmale noch Spielraum lassen, denn verbreitete Benchmarks bieten kaum noch welchen~\cite{zhang_revisiting_2025}. Die Vermutung aus der Einleitung, dass ein großer Teil des möglichen Beitrags ungelabelter Bilder bereits im Vortraining steckt, trifft damit eher auf \ac{CUB} und Semi-iNat-2021 zu als auf Cars196 und In-Shop bei kleinen Budgets. Schwache Merkmale allein sichern allerdings keinen großen Gewinn. Auf In-Shop sind die unveränderten Merkmale viel schwächer als auf Cars196, der größte Gewinn fällt aber kleiner aus (Abschnitt~\ref{sec:e:phase1}). Das könnte daran liegen, dass auf In-Shop beim kleinsten Budget schon rund \qty{31}{\percent} des Pools gelabelt sind, auf Cars196 nur \qty{4.9}{\percent} (Abschnitt~\ref{sec:budgets}). Zudem gewinnt die Referenz auf In-Shop bis zu allen Labels weniger hinzu (Tabelle~\ref{tab:e:referenz}).

Für \ac{CUB} liegt eine Erklärung in der hohen Ausgangsleistung auf den Validierungsfolds nahe (Abschnitt~\ref{sec:e:ausgangslage}). Wenn auf den Trainingsklassen ähnlich wenig zu ordnen bleibt, liefern die Labels dem Verlust kaum etwas, das die Merkmale nicht schon enthalten, was auch erklären würde, warum zusätzliche echte Labels die \UEBL dort kaum verbessern. Auf den Trainingsklassen selbst wurde das nicht gemessen. An der Auswahl zwischen den 100 Kombinationen liegt es vermutlich nicht, denn die Validierung trennt sie über dem Rauschen der Seeds und auch die im Test beste Kombination liegt kaum über den Merkmalen. Offen bleibt die Wahl der Epoche, weil auf \ac{CUB} nur das gewählte Modell getestet wird. Auf Semi-iNat-2021 spricht die ähnlich hohe \ac{P@1} vor dem Training für dieselbe Erklärung, die deutlich niedrigere Validierungs-\ac{MAP@R} dagegen.

\subsection{Label-Budget, Verlustfunktion und Aggregation}

Die zweite Forschungsfrage (F2) erfragte im ersten Teil, wie der Nutzen von dem Label-Budget, von dem Datensatz und von der Verlustfunktion abhängt. Vom Datensatz hängt er über den verbleibenden Spielraum ab, soweit vier Datensätze das erkennen lassen. Innerhalb eines Datensatzes sinkt der Gewinn der Verfahren, die bei kleinem Budget gewinnen, mit dem Budget. Der Gewinn von \ac{LRML} als Einzelmodell bleibt auf Cars196 etwa gleich (Abschnitt~\ref{sec:e:faktoren}). H2 ist damit als Muster teilweise gestützt. Das entspricht der Erwartung aus Abschnitt~\ref{sec:einordnung}, die sich auf Befunde der semi-supervised Klassifikation stützt~\cite{oliver_realistic_2018}. Frühere Arbeiten zum \ac{SSDML} erlauben hier dagegen keinen Vergleich, denn sie untersuchen je Datensatz nur ein Budget, wie fünf Labels je Klasse~\cite{dutta_semi-supervised_2021, zhuang_deep_2021, zhuang_deep_2023} oder alle Labels mit einem externen ungelabelten Datensatz~\cite{duan_slade_2021}. In dieser Arbeit bleibt die Ursache des Rückgangs offen, denn mit dem Budget ändern sich zugleich das Verhältnis ungelabelter zu gelabelten Bildern, die Stärke der Referenz und der Abstand zu dem Budget, bei dem die Hyperparameter gesucht wurden. Über einen Rückgang hinaus geht die Umkehr von \ac{LP} auf Semi-iNat-2021 bei $k=100$ (Abschnitt~\ref{sec:e:phase1}). Sie passt zum verbleibenden Spielraum, denn zusätzliche echte Labels bringen der Referenz dort nichts mehr (Tabelle~\ref{tab:e:referenz}), sodass selbst überwiegend richtige Pseudo-Labels vor allem durch ihre Fehler wirken könnten.%


Für das supervised \ac{DML} wurde gezeigt, dass sich Verlustfunktionen unter einheitlichen Versuchsbedingungen nur wenig unterscheiden~\cite{musgrave_metric_2020, roth_revisiting_2020}. Auch beim Fine-Tuning des ganzen DINOv2 liegen die Ergebnisse der Verlustfunktionen in einem Vergleich verschiedener Batchgrößen auf Cars196 und \ac{CUB} nahe beieinander~\cite{berton_all_2025}. Für den Nutzen eines Mechanismus kann die Verlustfunktion dagegen entscheidend sein. Auf Cars196, In-Shop und Semi-iNat-2021 liegt \ac{LP} bei kleinen Budgets als Einzelmodell mit jeder Verlustfunktion über der \UEBL, während bei \ac{LRML}, Hoffer und \ac{SLADE} schon das Vorzeichen des Gewinns von der Verlustfunktion abhängen kann (Abschnitt~\ref{sec:e:beste}). Dazu passt, dass \ac{LP} auf Cars196 bei jedem Budget mit \ac{MS} vorn liegt, der Verlustfunktion, die auch für die \UEBL meist am besten ist (Abschnitt~\ref{sec:e:beste}). Die Beiträge des Verfahrens und der Verlustfunktion scheinen sich dort zu addieren. Das passt dazu, dass \ac{LP} die Pseudo-Labels wie echte Labels an die Verlustfunktion übergibt, sodass kein zusätzlicher Term gegen den supervised Verlust abgewogen werden muss. Für die Bewertung eines Mechanismus mit eigenem Term genügt deshalb eine einzelne Verlustfunktion nicht, denn mit \ac{MS} verliert \ac{LRML} auf Cars196 bei $k=4$, mit den übrigen vier gewinnt es (Abbildung~\ref{fig:e:verluste}). Eine mögliche Erklärung liegt in der Gewichtung dieser Terme. Ihr Gewicht folgt einem Zielverhältnis zu einer Statistik der supervised Gradienten (Abschnitt~\ref{sec:m:gewicht}) und wie stark diese Gradienten auf wenige Parameter konzentriert sind, könnte von der Skala und der Temperatur der Verlustfunktion abhängen. Dasselbe $\eta$ würde dann je nach Verlustfunktion eine andere wirksame Stärke bedeuten. Dafür spricht, dass die Leistung von \ac{LRML} auf In-Shop mit Circle schon ab $\eta=\num{0.75}$ einbricht, mit Proxy-Anchor erst bei $\eta=3$ (Abschnitt~\ref{sec:e:sweeps}), gemessen wurde die wirksame Stärke aber nicht. Für die Praxis folgt daraus, dass sich die Konfiguration eines zusätzlichen Terms nicht ohne neue Suche von einer Verlustfunktion auf eine andere übertragen lässt. Eine Gewichtung, die das Verhältnis der euklidischen Gradientennormen laufend verfolgt, statt es nur zu Beginn zur Kalibrierung zu nutzen (Abschnitt~\ref{sec:m:gewicht}), könnte diese Abhängigkeit verringern und damit auch die Ergebnisse von \ac{LRML}, Hoffer und \ac{SLADE} verbessern.

Wie ein Gewinn zu bewerten ist, hängt außerdem davon ab, ob ein Einzelmodell oder ein Ensemble der Fold-Modelle eingesetzt wird. Bei \ac{LRML} auf Cars196 ist der Gewinn bei kleinen Budgets im Ensemble ein Mehrfaches des Gewinns als Einzelmodell, denn seine Fold-Modelle irren sich seltener gleich (Abschnitt~\ref{sec:e:phase2:concat}). Ein Mechanismus kann also vor allem die Vielfalt der Fold-Modelle erhöhen. Wer nur ein einzelnes Modell einsetzen kann, hat von diesem Vorteil nichts. Die Dimension der Embeddings im Index wächst mit der Zahl der Fold-Modelle, hier auf das Vierfache und auf \ac{CUB} auf das Fünffache (Abschnitt~\ref{sec:protokoll}).%

\section{Welche ungelabelten Daten ein Gewinn voraussetzt}
\label{sec:d:daten}

Neben dem Verfahren und der Verlustfunktion spielt auch die Zusammensetzung der ungelabelten Daten eine Rolle. In der In-Shop-Serie mit vollständig gelabelten Klassen stammen sie ausschließlich aus Klassen ohne Label. Externe Bildmengen wie in der Originalarbeit von \ac{SLADE} (Abschnitt~\ref{sec:ssdml}) prüft ein eigener Versuch (Abschnitt~\ref{sec:d:extern}).

\subsection{Bilder aus Klassen ohne Labels}

Die deutlichste Voraussetzung betrifft die Klassen. \ac{LP} und Hoffer ordnen jedes ungelabelte Bild einer gelabelten Klasse zu. Stammen die ungelabelten Bilder aus Klassen ohne Label, ist jede solche Zuordnung falsch und \ac{LP} ist in der In-Shop-Serie im Mittel über die Verlustfunktionen bei jedem Anteil gelabelter Klassen schlechter als die Referenz. Die Konfidenz erkennt Bilder neuer Klassen nicht, denn sie misst nur, wie stark sich die Score-Masse eines Bildes auf wenige bekannte Klassen konzentriert (Abschnitt~\ref{sec:lp}). Wie sehr die falschen Pseudo-Labels schaden, hängt aber von der Verlustfunktion ab. ArcFace und Proxy-Anchor, mit denen \ac{LP} am meisten verliert, ziehen jedes Bild zum Proxy seiner Klasse und damit Bilder neuer Klassen in eine fremde Klasse (Abschnitt~\ref{sec:verluste:proxy}). Paarbasierte Verlustfunktionen prüfen dagegen nur, ob zwei Bilder dasselbe Label tragen, sodass Bilder eines neuen Artikels mit demselben falschen Label trotzdem richtige positive Paare bilden können. Mit ArcFace, Proxy-Anchor und, gegen diese Erklärung, auch mit Circle erholt sich die Validierung nach den ersten Pseudo-Labels meist nicht mehr. Weil die Modellauswahl nach dem Warm-up neu beginnt (Abschnitt~\ref{sec:m:ssl}), kann \ac{LP} dann nicht auf das bessere Modell davor zurückfallen. H3 ist damit für \ac{LP} in der hier verwendeten Form gestützt. Ob eine absolute Konfidenzschwelle oder die mit der Konfidenz gewichtete Kreuzentropie des Originalverfahrens~\cite{iscen_label_2019} den Schaden begrenzen würde, wurde nicht untersucht. Für Hoffer ist H3 nur schwach gestützt, weil dessen Ergebnis über alle Anteile nahe null bleibt und nicht erkennbar ist, ob der Term dort wirkungslos ist oder den Bildern neuer Klassen standhält. Für \ac{SLADE}, das schon in der $k$-Serie auf In-Shop nicht gewinnt, lässt sich H3 hierdurch nicht prüfen.

Eine erwartete Ausnahme für Mechanismen ohne Zuordnung zu einer bekannten Klasse zeigt sich nur teilweise. \ac{LRML} gewinnt im Ensemble nur bei kleinen Anteilen gelabelter Klassen, mit Circle aber bei jedem Anteil (Abschnitt~\ref{sec:e:klassen}). Bei \qty{10}{\percent} gelabelter Klassen geht dieser Gewinn zum größten Teil auf die ungelabelten Bilder zurück (Tabelle~\ref{tab:e:klassen:2x2}). Ob dabei die Herkunft oder die Menge der ungelabelten Bilder entscheidet, trennt die Serie nicht, denn der ungelabelte Pool ist gerade bei den kleinen Anteilen am größten.

Für die semi-supervised Klassifikation wurde berichtet, dass Klassen, die nur im ungelabelten Pool vorkommen, die Leistung verschlechtern können~\cite{oliver_realistic_2018}. Dort wurden die Hyperparameter dafür nicht neu gesucht. Hier wurden die Hyperparameter dagegen eigens für die Serie mit \qty{10}{\percent} gelabelter Klassen gesucht (Abschnitt~\ref{sec:budgets}, Tabelle~\ref{tab:inshop10:hpo}). Die In-Shop-Serie dieser Arbeit deutet für das \ac{DML} an, dass der Schaden von der Art des Mechanismus abhängen könnte. Diese mögliche Abhängigkeit vom Mechanismus passt zu den kontrollierten Klassifikationsversuchen aus Abschnitt~\ref{sec:ssl:evaluation}. Die bereits beschriebenen Befunde zu Hoffer und \ac{SLADE} erlauben keine eindeutige Trennung nach dem Zuordnungsprinzip. Zudem beruht der Vergleich auf einem Datensatz und zwei Serien.

\subsection{Externe ungelabelte Bilder}
\label{sec:d:extern}

Für \ac{SLADE} wurde mit allen Labels und einem externen ungelabelten Datensatz ein deutlicher Gewinn gegenüber dem Teacher berichtet, auf Cars196 mit CompCars als externem Datensatz~\cite{duan_slade_2021}. In dieser Arbeit liegt unter ähnlichen Bedingungen kein Mechanismus über der \UEBL, auch die angepasste Variante von \ac{SLADE} nicht (Abschnitt~\ref{sec:e:extern}). Die Bedingungen unterscheiden sich aber, denn das Backbone ist hier eingefroren, die Referenz ist die \UEBL statt des Teachers (Abschnitt~\ref{sec:d:nutzen}) und die Cluster stammen vom trainierten Head selbst (Abschnitt~\ref{sec:m:slade}).

Für den Verlust kommen mehrere Gründe in Frage. Die Konfigurationen wurden für wenige Labels und den internen Pool gesucht und unverändert übernommen. Der Mechanismus kann bei vollständigen Labels auch ohne ungeeignete Bilder schaden. Die externen Bilder können zudem schaden, weil sie zu anderen Klassen gehören oder anders aussehen. Dass der Verlust bei \ac{LP}, Hoffer und \ac{SLADE} auf Cars196 schon mit \qty{10}{\percent} von CompCars voll auftritt und mit dem Pool kaum wächst, spricht zudem für einen Grund im Aufbau des Trainings. Bei \ac{LP} stammen auf Cars196 nur 12,5 bis \qty{25}{\percent} jedes Batches aus den gelabelten Bildern und der Rest aus pseudo-gelabelten Bildern, unabhängig von der Größe des Pools (Abschnitte~\ref{sec:m:lp} und~\ref{sec:e:extern}). Bei Hoffer und \ac{SLADE} legt das Zielverhältnis das Gewicht des Terms fest, ebenfalls unabhängig vom Pool (Abschnitt~\ref{sec:m:gewicht}). 

Bei \ac{LRML} läuft der Term in der Kontrolle weiter, auf Cars196 entsteht sein Leistungsverlust also erst durch die externen Bilder. Auf In-Shop liegen die Kontrollen im Mittel unter der \UEBL, sodass dort auch die übernommenen Konfigurationen als Grund in Frage kommen. Bei \ac{LP} gilt das aber nur für ArcFace. Mit den übrigen vier Verlustfunktionen liegt seine Kontrolle über der \UEBL und der Verlust entsteht erst mit den externen Bildern (Abschnitt~\ref{sec:e:extern}). Dazu passt die Art dieser Bilder. Da jede Klasse ein einzelner Artikel ist, dürfte kaum ein Bild aus Fashion200k zu einer gelabelten Klasse gehören, zudem sind die Bilder von Fashion200k anders als die von In-Shop zugeschnitten (Abschnitt~\ref{sec:daten}). Für \ac{LP} gleicht die Lage dann der In-Shop-Serie mit vollständig gelabelten Klassen, in der jedes Pseudo-Label falsch ist (Abschnitt~\ref{sec:d:daten}). Dass andere Nachbar- und Clusterzahlen die Verluste auf Cars196 nicht aufheben (Abschnitt~\ref{sec:e:extern}), spricht gegen diese beiden Werte als Ursache, schließt eine unpassende Konfiguration aber nicht aus, weil für dieses Setting keine vollständige \ac{HPO} durchgeführt wurde.%





\subsection{Nachbarschaften und Klassengröße}

Innerhalb der gelabelten Klassen hängt der Nutzen davon ab, wie gut die Nachbarschaften im Embedding-Raum des Heads die Klassen wiedergeben. Auffällig ist dabei, wie wenig reine Nachbarschaften \ac{LP} mit großen Graphen braucht. Bei $k=4$ liegt auf Cars196 nur mit 38 und 92 Nachbarn die Genauigkeit der ersten Pseudo-Labels über dem Anteil der Kanten innerhalb einer Klasse (Abschnitt~\ref{sec:e:verlauf}). Das spricht dafür, dass ein Pseudo-Label bei großen Graphen der Mehrheit vieler Pfade und nicht einzelnen Kanten folgt. Die Smoothness- und die Cluster-Annahme müssten für \ac{LP} dann nicht für jede Kante gelten, sondern nur für die Klassen als Ganzes. Für die offene Frage aus der Einleitung nach den Annahmen über die Struktur der Daten wäre das nur eine teilweise Antwort, denn sie gilt für die Trainingsklassen von Cars196 im Embedding-Raum nach dem Warm-up. Dass der Anteil gleichklassiger Kanten während des Trainings mit \ac{LP} steigt, ist allerdings kein unabhängiger Beleg dafür, dass diese Annahmen im Embedding-Raum zunehmend gelten, denn \ac{LP} trainiert den Head gerade mit den Labels, die aus diesem Graphen stammen. Wie sehr dabei die Ausgangsmerkmale zählen, deutet ein Vergleich an. Mit den Merkmalen eines vortrainierten BN-Inception und fünf Labels je Klasse ordnet Label Propagation auf ganz Cars196 nur rund ein Fünftel der Bilder richtig zu~\cite{zhuang_deep_2023}. Hier sind bei $k=4$ schon die ersten Pseudo-Labels mit ArcFace und Circle zu \qty{66}{\percent} bzw.\ \qty{78}{\percent} richtig (Abschnitt~\ref{sec:e:verlauf}). Diese Anteile beziehen sich allerdings nur auf die übernommenen Pseudo-Labels über der Konfidenzschwelle.

Für \ac{LRML} gilt diese Mehrheitsentscheidung nicht, denn der Term zieht jede gezogene Kante einzeln zusammen. Die beste Nachbarzahl könnte deshalb von der Klassengröße abhängen. Auf In-Shop mit wenigen Bildern je Klasse sind die kleinsten Nachbarzahlen am besten und reinere Kanten helfen deutlich, auf Cars196 mit vielen Bildern je Klasse liegen viele Nachbarzahlen gleichauf und reinere Kanten helfen kaum (Abschnitt~\ref{sec:e:sweeps}, Tabelle~\ref{tab:e:ablationen:verluste}). Bei wenigen Bildern je Klasse verlassen schon wenige Nachbarn die Klasse und der Term erzwingt Nähe genau dort, wo die Labels wechseln. Die Zahl der Bilder je Klasse im Pool könnte damit eine Größe sein, nach der sich die Nachbarzahl ohne Labels wählen lässt, sofern die ungefähre Klassengröße bekannt ist. Auf Cars196 wirkt der Term dagegen auf andere Weise (Abschnitt~\ref{sec:d:ablationen}). Dieselbe Schwierigkeit ist für ein graphbasiertes Verfahren beschrieben, das negative Paare über eine Nachbarzahl bestimmt, besonders bei unausgeglichenen Klassen~\cite{zhuang_deep_2023}. In dieser Arbeit gewinnt \ac{LRML} auf Semi-iNat-2021 nicht, dessen Arten im Mittel rund 150, aber sehr unterschiedlich viele Bilder haben (Abschnitt~\ref{sec:budgets}). Dort könnte keine einzelne Nachbarzahl zu großen und kleinen Arten zugleich passen.

\section{Was die Ablationen erklären}
\label{sec:d:ablationen}

Für F2 sollte außerdem geklärt werden, worauf ein Gewinn beruht. Die separat abgestimmten Hyperparameter allein erzeugen ihn nicht (Abschnitt~\ref{sec:e:2x2}). Die Kontrolle mit abgeschaltetem \ac{SSL}-Term liegt in 25 der 40 Kombinationen sogar um mehr als die Rauschschwelle unter der \UEBL, nicht wie erwartet auf ihrem Niveau. Das spricht dafür, dass die Suche die Hyperparameter auf den jeweiligen Mechanismus abstimmt. Der Vergleich mit dieser Kontrolle zeigt deshalb nur, was der Term innerhalb seiner Konfiguration bewirkt, nicht schon einen Gewinn gegenüber der \UEBL. Wie unsicher ein Urteil mit einem Seed ist, zeigt die Zerlegung selbst, denn sechs der 14 Kombinationen ohne Gewinn liegen im Mittel der fünf Seeds der Phase~1 über der \UEBL (Abschnitt~\ref{sec:e:2x2}).

\subsection{Ursachen je Mechanismus}

Für \ac{LP} ist die Zerlegung am klarsten. Der Mechanismus bewirkt allein auf den gelabelten Bildern fast nie etwas und sein Gewinn bleibt mit den gemeinsamen Hyperparametern der Referenz meist bestehen (Abschnitt~\ref{sec:e:2x2}). Er beruht also auf den ungelabelten Bildern und, wie der umverdrahtete Graph zeigt (Abschnitt~\ref{sec:e:nachbarschaft}), auf ihrer Nachbarschaft. Dass es auf die aus den Daten berechnete Nachbarschaft ankommt, passt zur Klassifikation, wo der Vergleich von Pseudo-Labels aus dem Graphen und aus den Vorhersagen des Netzes zugunsten des Graphen ausfällt~\cite{iscen_label_2019}.%


Bei \ac{LRML} ist der Gewinn auf Cars196 nur zum Teil ein Gewinn durch ungelabelte Bilder. Schon auf den gelabelten Bildern hebt der Term die Leistung mit allen fünf Verlustfunktionen über die Kontrolle mit abgeschaltetem Term. Klassenreine Nachbarschaften erklären diesen Beitrag nicht. Die in Abschnitt~\ref{sec:e:2x2} hergeleitete Mehrheit klassenübergreifender Nachbarschaftsbeziehungen in dieser Variante spricht gegen eine Erklärung allein durch klassenreine Kanten. Zudem helfen reinere Kanten \ac{LRML} auf Cars196 kaum. Ein zufälliger Graph schadet dagegen. Das spricht dafür, dass der Term dort weniger die Klassenstruktur überträgt, als vielmehr als datenabhängige Glättung wirkt, die benachbarte Merkmale ähnlich abbildet. Das entspricht der Idee der Mannigfaltigkeitsregularisierung, die mit einem datenabhängigen Term Funktionen bevorzugt, die entlang der Datenverteilung glatt sind~\cite{belkin_manifold_2006}. Auf In-Shop und im Ensemble tragen dagegen vor allem die ungelabelten Bilder bei (Abschnitte~\ref{sec:e:2x2} und~\ref{sec:e:phase2:concat}).%


Bei Hoffer stammt der Gewinn vor allem aus den ungelabelten Bildern. Auf In-Shop hebt aber auch der Term auf den gelabelten Bildern allein die Leistung. Mit den gemeinsamen Hyperparametern der Referenz gewinnt Hoffer auf In-Shop, auf Cars196 nicht (Abschnitt~\ref{sec:e:2x2}). Ob das Ausbleiben auf Cars196 eine echte Abhängigkeit von der übrigen Konfiguration anzeigt oder nur eine unpassende Kombination übertragener Werte, trennen die Versuche nicht.%


Bei \ac{SLADE} zeigen die Ablationen, dass die ungelabelten Bilder auf In-Shop kaum etwas beitragen (Abschnitt~\ref{sec:e:2x2}), aber nicht, warum. Weder die eigenen noch die Hyperparameter der \UEBL führen dort zu einem Gewinn und kein Wert der Sweeps verbessert das Ergebnis deutlich (Abschnitte~\ref{sec:e:2x2} und~\ref{sec:e:sweeps}). Andere gemeinsame Einstellungen mehrerer Hyperparameter sind damit nicht ausgeschlossen. Eine mögliche Erklärung liegt in den sicheren Paaren, die \ac{SLADE} aus Clustern ohne Bezug zu den gelabelten Klassen bildet. Bei wenigen Bildern je Artikel könnten viele dieser Paare verschiedene Artikel verbinden, so wie bei \ac{LRML} auf In-Shop schon wenige Nachbarn die Klasse verlassen (Abschnitt~\ref{sec:d:daten}). Auch im Original erreicht \ac{SLADE} auf In-Shop nur das Niveau supervised trainierter Verfahren~\cite{duan_slade_2021}. Auf Cars196 gewinnt die hier verwendete Form ohne getrennten Teacher dagegen bei kleinen Budgets.%


\subsection{Falsche Pseudo-Labels}

Die Genauigkeit der Pseudo-Labels allein ist ein unzureichendes Maß für ihren Nutzen. Bei ähnlich genauen ersten Pseudo-Labels gewinnt \ac{LP} bei $k=2$ auf In-Shop ein Mehrfaches des Gewinns auf Semi-iNat-2021 (Abschnitt~\ref{sec:e:pseudolabels}). Zum Teil passt das dazu, dass echte Labels die Referenz von $k=2$ bis zu allen Labels auf Semi-iNat-2021 nur etwa halb so stark verbessern wie auf In-Shop (Tabelle~\ref{tab:e:referenz}). Der Unterschied der Gewinne ist aber deutlich größer. Auf Semi-iNat-2021 liegt \ac{LP} bei hundert Labels je Art sogar trotz überwiegend richtiger Pseudo-Labels in allen Vergleichen unter der Referenz (Abschnitt~\ref{sec:e:pseudolabels}). Gegenüber den Oracle-Varianten liegt reguläres \ac{LP} meist über der Kurve zufälliger Fehler gleicher Häufigkeit (Abschnitt~\ref{sec:e:pseudolabels}). Der Vergleich verwendet dieselbe mittlere Fehlerrate, bildet aber weder deren zeitlichen Verlauf noch die Auswahl der Bilder zwischen den Varianten identisch ab. Die veränderten Labels können auch spätere Graphen und Auswahlschritte beeinflussen. Der Befund beruht zudem auf einem Seed. Die kleine Wirkung zufälliger Fehler mit ArcFace und \ac{NT-Xent} auf Cars196 passt zum Befund an verrauschten Trainingslabels aus Abschnitt~\ref{sec:m:phase2}. Dass Circle, \ac{MS} und Proxy-Anchor hier ihren Gewinn schon bei \qty{20}{\percent} verlieren, widerspricht dem nicht unbedingt, denn ein Gewinn von wenigen Prozentpunkten verschwindet schon bei einem Schaden, der gemessen an der ganzen Leistung gering ist. Auf In-Shop schaden zufällige Fehler dagegen mit jeder Verlustfunktion deutlich (Abschnitt~\ref{sec:e:pseudolabels}), vermutlich weil ein falsches Label bei wenigen Bildern je Artikel einen größeren Teil einer Klasse betrifft. Die Fehler von \ac{LP} sind strukturiert, die meisten falschen Pseudo-Labels nennen ein Modell desselben Herstellers. Auch in der Klassifikation konzentrieren sich die sicheren Fehler von \ac{LP} auf wenige verwechselbare Klassenpaare wie Lastwagen und Autos~\cite{iscen_label_2019}. Ob diese Struktur in dieser Arbeit den Abstand zur Kurve verursacht, zeigt der Vergleich aber nicht, denn zufällige Fehler gleicher Häufigkeit ändern zugleich, welche Bilder betroffen sind und welche Klassen verwechselt werden. Entscheiden würde eine Oracle-Variante, deren falsche Labels dem Verwechslungsmuster von \ac{LP} folgen. Falsche Pseudo-Labels sind also nicht harmlos, richtige Labels an denselben Stellen würden \ac{LP} weiter verbessern. Reguläres \ac{LP} liegt aber über den Läufen mit zufälligen Fehlern gleicher Häufigkeit.%


\subsection{Geometrie und vorab festgelegte Erwartungen}

Die vorab formulierte Erwartung, dass die Referenz bei wenigen Labels stärker komprimiert als die \ac{SSL}-Heads, trifft auf keinem der drei untersuchten Datensätze zu. Für vollständig gelabelte Daten und ein mittrainiertes Backbone wurde berichtet, dass weniger komprimierte Räume besser auf ungesehene Klassen generalisieren und auf \ac{CUB} und Cars196 ein größeres $\pi_{\mathrm{ratio}}$ schwächer mit besserer Generalisierung einhergeht~\cite{roth_revisiting_2020}. Auf Cars196 gehen in dieser Arbeit dagegen kompaktere Klassen mit besserem Retrieval einher. Bei den kleinen Budgets lässt sich das aber nicht von der Dauer des Trainings bis zur gewählten Epoche trennen. Der Zusammenhang gilt aber nicht für jedes Verfahren, denn \ac{LRML} erzeugt nach \ac{LP} die kompaktesten Klassen, erreicht in den Geometrieläufen bei $k=4$ als Einzelmodell aber nur das Niveau der Referenz (Abschnitt~\ref{sec:e:geometrie}), was zur Deutung des Terms als Glättung passt. Auf Semi-iNat-2021 und In-Shop ist der Zusammenhang mit $\pi_{\mathrm{ratio}}$ innerhalb der kleinen Budgets negativ und innerhalb der größeren positiv (Tabelle~\ref{tab:e:korrelation}), sodass der negative Zusammenhang über alle Budgets zum Teil aus den Unterschieden zwischen den Budgets entsteht. Als Kennzahl für die Wahl zwischen Konfigurationen desselben Budgets kommt $\pi_{\mathrm{ratio}}$ hier deshalb allenfalls bei kleinen Budgets in Frage. Seine Berechnung für den gesamten Trainingspool setzt jedoch die verdeckten Klassenlabels voraus (Abschnitt~\ref{sec:m:geometrie}). Eine Modellauswahl ohne diese Labels ist damit nicht belegt. Da der spektrale Abfall schon durch die zufällige Projektion des Heads steigt, ohne dass sich die Leistung ändert, beschreibt er hier eher die Parametrisierung des Heads als die Generalisierung. Über die ooDML-Aufteilungen von \ac{CUB} und Stanford Online Products hinweg zeigten $\pi_{\mathrm{ratio}}$ und der spektrale Abfall ebenfalls kaum zuverlässig an, welches Modell innerhalb einer Aufteilung die höhere Recall@1 erreicht~\cite{milbich_characterizing_2021}. In dieser Arbeit scheitert die Erwartung damit eher an ihrem Indikator, dem spektralen Abfall, als an der Überanpassung selbst, die sich auf Cars196 bei $k=4$ an der wachsenden Streuung außerhalb der gelabelten Bilder zeigt (Abschnitt~\ref{sec:e:verlauf}). Die enge Korrelation der Genauigkeit der Prototypen-Zuordnung mit der Test-\ac{MAP@R} ist hier zum Teil durch das Training vorgegeben, weil die ungelabelten Bilder für die Mechanismen Trainingsdaten sind. Dass unter \ac{LP} anders als unter der Referenz auch die ungelabelten Bilder und die Validierungsklassen kompakter werden (Abschnitt~\ref{sec:e:verlauf}), passt zu der zusätzlichen Annahme aus Abschnitt~\ref{sec:ssl}, dass sich das an den ungelabelten Bildern Gelernte auf ungesehene Klassen überträgt.%


Nur drei der vorab festgelegten Erwartungen der Phase~2 wurden vollständig bestätigt (Abschnitt~\ref{sec:e:sweeps}). Dass mehr Nachbarn \ac{LP} auf Cars196 meist wenig ausmachen, auf In-Shop mit wenigen Bildern je Klasse aber teils schaden, passt zur Mehrheitsentscheidung aus Abschnitt~\ref{sec:d:daten}, die nur trägt, solange genügend Pfade innerhalb einer Klasse verlaufen. Dass \ac{LP} meist kaum vom Warm-up abhängt, passt dazu, dass sein Graph hier von Beginn an auf vortrainierten Merkmalen entsteht. Im Originalverfahren entsteht der erste Graph dagegen aus einem Netz, das von zufälligen Gewichten aus nur auf den gelabelten Bildern trainiert wurde~\cite{iscen_label_2019}. Belastbar sind vor allem Effekte, die über Verlustfunktionen und Datensätze hinweg dieselbe Richtung haben, wie der Schaden eines umverdrahteten Graphen, der Beitrag der ungelabelten Bilder bei \ac{LP} und Hoffer und die Grenzen von Distanzskala und $\eta$.

\section{Aufwand und Hinweise für das Training}
\label{sec:d:praxis}

\subsection{Entscheidende Hyperparameter und Aufwand}

Die dritte Forschungsfrage (F3) lautete, welche Hyperparameter und Designentscheidungen bestimmen, ob ein Mechanismus hilft. Die Suchen haben flache Optima, sodass nicht die genauen Werte zählen, sondern wenige Grenzen, jenseits derer ein Mechanismus versagt (Abschnitte~\ref{sec:e:sweeps} und~\ref{sec:e:sensitivitaet}). Bei Hoffer sind es ein Mindestwert von $\eta$ und eine untere Grenze der Distanzskala, die auf In-Shop mit seinen viel zahlreicheren Klassen höher liegt als auf Cars196. Dass die nötige Skala bei mehr Klassen größer ist, passt zu der Schranke der Zuordnungswahrscheinlichkeit (Abschnitt~\ref{sec:m:hoffer}) und dazu, dass beim normierten Softmax-Verlust die untere Schranke des Verlusts bei fester Skala mit der Zahl der Klassen steigt~\cite{wang_normface_2017}. Bei \ac{LRML} ist es ein Fenster des Zielverhältnisses $\eta$, dessen obere Grenze von der Verlustfunktion abhängt, dazu eine je nach Datensatz passende Nachbarzahl und ein Warm-up, ohne das \ac{LRML} vor allem mit ArcFace und \ac{NT-Xent} einbricht. Der Einbruch bei großem $\eta$ passt dazu, dass ein Laplace-Term ohne abstoßenden Teil die Embeddings zusammenziehen kann (Abschnitt~\ref{sec:m:lrml}). Darauf deutet auch, dass unter \ac{LRML} auf Cars196 bei $k=4$ schon mit dem gewählten $\eta$ die Klassenmittelpunkte nach dem Warm-up mit vier der fünf Verlustfunktionen wieder zusammenrücken (Abschnitt~\ref{sec:e:verlauf}). Weil dort auch nach dem Training ein großer Teil der Kanten zwischen verschiedenen Klassen verläuft (Tabelle~\ref{tab:a:reinheit:lrml}), zieht der Term vermutlich auch benachbarte Klassen zueinander. Bei \ac{SLADE} schadet ebenfalls ein großes $\eta$, kein untersuchter Wert verbessert es aber deutlich (Abschnitt~\ref{sec:e:sweeps}). \ac{LP} braucht vor allem zwei Designentscheidungen, die Neuberechnung des Graphen und der Pseudo-Labels nach jeder Epoche und eine zur Verlustfunktion und zum Datensatz passende Konfidenzschwelle. Mit ArcFace und \ac{MS} auf Cars196 und mit Circle auf In-Shop darf sie nicht zu viele Bilder ausschließen, mit ArcFace auf In-Shop übernimmt \ac{LP} mit der gewählten Schwelle dagegen anfangs nur wenige Bilder (Abschnitt~\ref{sec:e:sweeps}). Die Neuberechnung je Epoche stammt aus dem Originalverfahren, in dem die Genauigkeit der Pseudo-Labels über das Training steigt~\cite{iscen_label_2019}. Gegenüber der Änderung einzelner übriger Hyperparameter ist \ac{LP} hier je nach Verlustfunktion und Datensatz robust. Am wenigsten robust ist ArcFace auf In-Shop (Abschnitt~\ref{sec:e:sweeps}). Die Neuberechnung ist damit insgesamt nützlich, begrenzt aber nicht nachweislich den Confirmation Bias (Abschnitt~\ref{sec:teacher-student}). Steigende Abdeckung oder Genauigkeit schließen nicht aus, dass einzelne falsche Pseudo-Labels bestehen bleiben. Dass im gezeigten Graphlauf auf Cars196 die meisten falschen Pseudo-Labels ihre falsche Klasse behalten (Abschnitt~\ref{sec:e:pseudolabels}), passt dazu, dass sich beim Confirmation Bias falsche Pseudo-Labels verfestigen~\cite{arazo_2020}. Ebenso kann es daran liegen, dass diese Bilder schon zu Beginn von \ac{LP} nahe bei der falschen Klasse liegen. Dass die Aktualisierung auf Cars196 trotzdem hilft, passt dazu, dass sie dort vor allem neue und überwiegend richtige Pseudo-Labels hinzufügt.

Mit gespeicherten Merkmalen spricht die Rechenzeit gegen keinen der Mechanismen, die helfen, denn auf In-Shop bei $k=2$ braucht auch der langsamste je Fold im Median nur gut zehn Minuten (Abschnitt~\ref{sec:e:aufwand}). Beim Fine-Tuning des Backbones müssten der Graph, die Pseudo-Labels oder die Cluster dagegen in jeder Epoche aus neu berechneten Embeddings des ganzen Pools entstehen.


\subsection{Bedingte Empfehlung}

Zusammengenommen ergibt sich eine bedingte Empfehlung. Wie die Versuche setzt sie eine gelabelte Validierung voraus, die groß genug für eine zuverlässige Modellauswahl ist. Ob sich ein Head überhaupt lohnt, zeigt erst der Vergleich mit den unveränderten Merkmalen auf diesen Validierungsklassen, denn auf \ac{CUB} übertrifft ein trainierter Head sie als Einzelmodell meist nicht. Die Größe eines Gewinns gegenüber einer Referenz lässt sich von den Validierungsklassen nicht ohne Weiteres auf den Test übertragen, denn auf Cars196 überschätzt die Validierung die Gewinne, am stärksten bei \ac{LP} (Abschnitt~\ref{sec:e:faktoren}). Wenn zusätzliche Labels die Referenz noch deutlich verbessern, nur wenige Bilder je Klasse gelabelt sind und die ungelabelten Bilder aus den gelabelten Klassen stammen, ist \ac{LP} die am besten gestützte Wahl. Unter diesen Bedingungen ist es auf Cars196 und In-Shop in keinem gepaarten Vergleich schlechter als die Referenz (Abschnitt~\ref{sec:e:phase1}). Als Einzelmodell gewinnen dann auch Hoffer und auf In-Shop \ac{LRML}, Hoffer aber nur mit einer ausreichend großen, empirisch abgestimmten Distanzskala und auf Cars196 nur mit einer gemeinsam abgestimmten Konfiguration. Ist ein großer Teil des Pools bereits gelabelt oder bringen zusätzliche Labels der Referenz wenig, wie auf Semi-iNat-2021 bei $k=100$, bleiben die Gewinne klein oder kehren sich um und die Referenz genügt. Als Verlustfunktion lag \ac{MS} in den meisten Fällen vorn, auch für die \UEBL, Proxy-Anchor war auf In-Shop am schwächsten (Abschnitt~\ref{sec:e:beste}). Wird das Ensemble der Fold-Modelle eingesetzt, fallen die Gewinne bei kleinen Budgets meist größer aus, vor allem bei \ac{LRML}. Auf In-Shop sind sie bei $k=2$ im Ensemble dagegen für drei der vier Verfahren kleiner als bei den Einzelmodellen (Abbildung~\ref{fig:e:gewinne}).

Stammen die ungelabelten Bilder aus Klassen ohne Label, etwa von neu aufgenommenen Artikeln eines Online-Shops, stützen die Ergebnisse keinen Mechanismus stark. Dort mischen sich zudem bekannte und neue Artikel, ein Fall, den die Versuche nicht abdecken (Abschnitt~\ref{sec:d:grenzen}). Für Bilder neuer Klassen kommt am ehesten \ac{LRML} in Frage, das bei kleinen Anteilen gelabelter Klassen im Ensemble und mit Circle auch als Einzelmodell gewinnt, wobei diese Gewinne mit dem Budget der Hyperparametersuche und dem größten ungelabelten Pool zusammenfallen (Abschnitt~\ref{sec:d:daten}). Wenn annotiert werden kann, spricht der Vergleich bei ähnlichem gelabeltem Anteil (Abschnitt~\ref{sec:e:klassen}) eher dafür, wenige Bilder jeder Klasse zu labeln als einzelne Klassen vollständig, auch wenn er zwei Serien mit verschiedenen Einstellungen verbindet. Gegenüber nicht verglichenen Alternativen wie einem mittrainierten Backbone, einer gezielten Auswahl der zu labelnden Bilder oder einer Nutzung des Nachbarschaftsgraphen erst bei der Suche liefert die Arbeit keine Rangfolge.

\section{Grenzen und zukünftige Arbeiten}
\label{sec:d:grenzen}

Die Reichweite der Aussagen ist zuerst durch den Aufbau begrenzt. Trainiert wurde ein linearer Head über den Merkmalen eines eingefrorenen Backbones (Abschnitt~\ref{sec:aufbau}). Beim Fine-Tuning des Backbones verändern ungelabelte Bilder auch die Merkmale selbst und Befunde wie die Wirkung der Laplace-Regularisierung als Glättung des Heads oder die Robustheit von \ac{LP}, die schon in diesem Aufbau von der Verlustfunktion und dem Datensatz abhängt, müssen dort nicht gelten. Naheliegend ist deshalb, die Vergleiche mit Fine-Tuning zu wiederholen. Einen Anhaltspunkt, was der eingefrorene Backbone an Leistung kostet, gibt ein Vergleich in einem anderen Aufbau auf \ac{CUB} mit einem auf ImageNet-21k vortrainierten ViT-S/16 und Proxy-Anchor. Wird dort nur ein linearer Head trainiert, liegt die \ac{MAP@R} drei Prozentpunkte unter dem Fine-Tuning des ganzen Netzes. Einige Verfahren, die nur einen kleinen Teil der Parameter im Backbone anpassen, etwa nur die Bias-Terme, erreichen dagegen die Werte des vollständigen Fine-Tunings oder liegen leicht darüber~\cite{ren_learning_2024}. Offen ist auch, ob eine gezielte Begrenzung der Kompression bei eingefrorenem Backbone hilft. Der spektrale Abfall aus Abschnitt~\ref{sec:geometrie} wäre dafür hier allerdings kein ausreichender Erfolgsindikator, weil er eher die Parametrisierung des Heads beschreibt (Abschnitt~\ref{sec:d:ablationen}). Mit den gespeicherten Merkmalen entfiel zudem die Konsistenzregularisierung. \ac{LRML} lässt sich als eine Form davon verstehen, deren zweite Ansicht eines Bildes sein Nachbar ist. Anders als eine Augmentierung teilt ein Nachbar das Label aber nur, soweit der Embedding-Raum die Klassen bereits trennt. Eine Konsistenz über mehrere Ansichten eines Bildes wäre auch ohne Training des Backbones möglich. Dafür müssten die Ansichten wie üblich in jedem Batch neu augmentiert und durch das eingefrorene DINOv2 geführt oder vorab mehrere augmentierte Ansichten je Bild gespeichert werden. In der Klassifikation verbesserte die Kombination von Label Propagation und Mean Teacher die Ergebnisse auf CIFAR-10 sowie mit \num{10000} Labels auf CIFAR-100 und Mini-ImageNet, auf Mini-ImageNet mit \num{4000} Labels war Label Propagation allein besser~\cite{iscen_label_2019}. Ob das auch im \ac{DML} gilt, ließe sich in einem solchen Aufbau prüfen.%

Für Cars196 und \ac{CUB} sind die Testklassen für DINOv2 wegen der Auswahl seiner Vortrainingsdaten nicht im strengen Sinn unbekannt (Abschnitt~\ref{sec:vit}). Die Gewinne auf In-Shop und Semi-iNat-2021 zeigen, dass ein Nutzen nicht auf die beiden für die Vortrainingsauswahl verwendeten Datensätze beschränkt ist (Abschnitt~\ref{sec:vit}). Welchen Anteil diese Auswahl an den Unterschieden zwischen den Datensätzen hat, lässt sich daraus nicht bestimmen. Dafür fehlt ein Vergleich desselben Aufbaus mit gezielt veränderten Vortrainingsdaten. Die Testklassen blieben hier je Datensatz fest. Auf den schwierigeren ooDML-Aufteilungen von Stanford Online Products verschiebt sich der Recall@1-Vergleich zuungunsten von ArcFace und Proxy-Anchor gegenüber den übrigen Verfahren~\cite{milbich_characterizing_2021}. Für diese Arbeit bleibt offen, ob die Gewinne von \ac{LP} und die geringen Unterschiede zwischen den Referenzen der fünf Verlustfunktionen auf anderen Klassenaufteilungen bestehen bleiben. Nicht untersucht sind außerdem ungelabelte Bilder aus einer deutlich anderen Domäne, gemischte Pools aus bekannten und neuen Klassen bei wenigen Labels, Duplikate sowie mehrdeutige oder verrauschte Bilder. Einen solchen gemischten Pool ergäben die hier nicht verwendeten Bilder der restlichen Arten ohne Labels aus Semi-iNat-2021 zusammen mit den ungelabelten Bildern der Trainingsarten (Abschnitt~\ref{sec:daten}). Nach den Befunden aus Abschnitt~\ref{sec:d:daten} wäre für \ac{LP} mit einem solchen Pool eher ein kleinerer Gewinn als in der $k$-Serie oder ein Verlust zu erwarten. Bei vollständigen Labels dürften Bilder einer deutlich anderen Domäne kaum helfen, denn schon die externen Bildmengen CompCars und Fashion200k bringen dort keinem Mechanismus im Mittel einen Gewinn (Abschnitt~\ref{sec:d:extern}).%

Der in Abschnitt~\ref{sec:budgets} beschriebene Umfang der vollständig gelabelten Validierung begrenzt die Übertragbarkeit auf ein Setting mit insgesamt knappen Labels. Ob sich Konfigurationen und Modelle finden lassen, die mit deutlich weniger gelabelten Validierungsdaten gewählt werden und ihren Vorsprung behalten, ist offen. Das betrifft besonders Mechanismen mit vielen Hyperparametern. Wie die unterschiedlichen Suchumfänge aus Abschnitt~\ref{sec:m:phase1} unter knapper Validierung zu bewerten sind, bleibt damit offen. Dieselbe Zahl von 80 Versuchen je Suche ist zudem nur eine mögliche Vorgabe. Bei gleicher Rechenzeit hätte die schnellere Referenz mehr Versuche erhalten und bei gleicher Dichte der Suche die Mechanismen mit mehr Hyperparametern. Die flachen Optima beschreiben die bereits untersuchten Bereiche (Abschnitt~\ref{sec:e:sensitivitaet}). Ob eine andere Vorgabe die Rangfolge der Verfahren verändert, bleibt offen.%

 Für Semi-iNat-2021 waren Ablationen wie das Hinzufügen der restlichen Klassen als ungelabelte Daten geplant, wurden aus Zeitgründen aber nicht durchgeführt, sodass die Vermutung zur Nachbarzahl von \ac{LRML} (Abschnitt~\ref{sec:d:daten}) auf zwei Datensätzen beruht. In der In-Shop-Serie fehlt ebenso ein Versuch, der Herkunft und Menge der ungelabelten Bilder trennt, etwa mit einem verkleinerten Pool bei gleichem Anteil gelabelter Klassen. Für die in Phase~2 nicht gezielt variierte Zahl der Hoffer-Referenzmengen (Abschnitt~\ref{sec:m:hoffer}) fehlt ebenfalls ein isolierter Vergleich. Ob mehrere Referenzmengen die Zuordnung der ungelabelten Bilder stabilisieren, bleibt deshalb offen.%

Am wenigsten erklärt ist das Abschneiden von \ac{SLADE} auf In-Shop. Gemessen wurden weder die Reinheit der Cluster noch die Genauigkeit der sicheren Paare. Zudem wurde nur eine angepasste Form ohne getrennten Teacher untersucht (Abschnitt~\ref{sec:m:slade}). Eine Auswertung wie für die Pseudo-Labels von \ac{LP} (Abschnitt~\ref{sec:e:pseudolabels}) wäre dafür der nächste Schritt. Offen bleibt auch, ob die Weitergabe des Teachers über die Folds zu den mit externen Bildern berichteten Gewinnen beiträgt, denn sie wurde wegen der Validierungsklassen nicht übernommen (Abschnitt~\ref{sec:m:slade}).%

Auch \ac{LRML} wurde nur mit einem ungewichteten Graphen und einem rein anziehenden Term untersucht (Abschnitt~\ref{sec:m:lrml}). Weil reinere Kanten \ac{LRML} auf In-Shop mit vier der fünf Verlustfunktionen verbessern (Tabelle~\ref{tab:e:ablationen:verluste}), liegt es nahe, die Variante zu prüfen, die den Laplace-Term nach Dichte und Ähnlichkeit gewichtet~\cite{ying_manifold_2018}. Einfacher wäre eine Gewichtung allein nach dem Abstand, bei der das Gewicht einer Kante exponentiell mit dem quadrierten Abstand ihrer Bilder fällt~\cite{belkin_manifold_2006}. Einen Schritt weiter ginge ein Graphterm, der Kanten als positive Paare und Paare weit entfernter Bilder als negative an eine \ac{DML}-Verlustfunktion wie \ac{MS} übergibt. Er würde der Embedding-Regularisierung (Abschnitt~\ref{sec:weston}) ähneln, aber statt des Contrastive Loss eine der hier untersuchten Verlustfunktionen verwenden. Der Miner von \ac{MS} könnte weit entfernte Negative allerdings vollständig verwerfen (Abschnitt~\ref{sec:verluste:paare}). Eine Variante der Label Propagation wählt dagegen gerade nahe Paare als Negative, indem sie geschätzte negative Kanten in den Graphen aufnimmt, um nahe Bilder verschiedener Klassen zu trennen~\cite{zhuang_deep_2023}.%

\section{Fazit}
\label{sec:d:reichweite}

Ungelabelte Bilder verbessern das \ac{DML} auf eingefrorenen DINOv2-Merkmalen gegenüber einer abgestimmten supervised Referenz nur unter bestimmten Bedingungen (F1). Am größten ist ihr Nutzen bei wenigen Labels und großem verbleibendem Spielraum. Er hängt außerdem von der Verlustfunktion ab und setzt bei \ac{LP} ungelabelte Bilder aus den gelabelten Klassen voraus (F2). Breit hilft nur \ac{LP}, dessen Gewinn auf den ungelabelten Bildern und ihrer Nachbarschaft beruht. Entscheidend sind wenige Grenzen der Hyperparameter und nicht ihre genauen Werte, bei \ac{LP} zudem die Neuberechnung des Graphen und der Pseudo-Labels (F3). Über die einzelnen Retrieval-Werte hinaus zeigt die Arbeit, wie der Nutzen ungelabelter Bilder im \ac{DML} bewertet werden sollte. Eine gleich intensiv abgestimmte supervised Referenz, wie sie für die semi-supervised Klassifikation empfohlen wird~\cite{oliver_realistic_2018}, ist notwendig, aber nicht ausreichend. Hinzu kommt in dieser Arbeit die Zerlegung nach den Hyperparametern, dem Term und den ungelabelten Bildern. Ohne sie wäre bei \ac{LRML} auf Cars196 der Beitrag der ungelabelten Bilder überschätzt worden. Weitere Beiträge sind die unveränderten Merkmale als eigener Bezugspunkt, getrennte Angaben für das Einzelmodell und das Ensemble und vorab festgelegte Erwartungen. Die Genauigkeit der Pseudo-Labels reicht als Kennzahl dagegen nicht aus. Für den Entwurf von Mechanismen spricht die Arbeit dafür, ungelabelte Bilder über Pseudo-Labels in die vorhandene Verlustfunktion einzubringen, sofern die ungelabelten Bilder aus den gelabelten Klassen stammen (Abschnitt~\ref{sec:d:daten}).








